\section{长期记忆（Long-term Memory）}

\subsection{短期记忆的局限}

LLM Agent 的``短期记忆''即上下文窗口（context window），存在三个根本限制：

\begin{enumerate}
    \item \textbf{只能追加}（append-only）：只能不断添加新 token，不能编辑或删除
    \item \textbf{容量有限}：即使有百万 token 窗口，仍有上限
    \item \textbf{不持久}：任务结束或上下文清空后，所有"记忆"消失
\end{enumerate}

\begin{warningbox}{金鱼困境}
没有长期记忆的 Agent 就像金鱼——即使今天证明了黎曼猜想，明天也得从头再来。更严重的是，下次尝试不一定能成功。
\end{warningbox}

\subsection{Reflexion：最简单的长期记忆}

Reflexion 是 ReAct 的自然扩展：

\begin{enumerate}
    \item Agent 尝试解决任务（如编写代码）
    \item 任务失败（如单元测试未通过）
    \item Agent \textbf{反思}失败原因（如 ``我忘记处理边界情况''）
    \item 将反思结果写入长期记忆
    \item 下次尝试时读取长期记忆，避免重复错误
\end{enumerate}

\begin{knowledgebox}{一种新型学习范式}
Reflexion 本质上是一种\textbf{通过语言进行学习}的方法，区别于传统 RL 通过梯度下降更新权重：
\begin{itemize}
    \item 不需要标量奖励信号，可以利用任何文本反馈（编译错误、测试结果、教师评语等）
    \item 不通过反向传播更新参数，而是通过更新文本记忆影响未来行为
\end{itemize}
\end{knowledgebox}

\subsection{更复杂的长期记忆形式}

\begin{itemize}
    \item \textbf{Voyager}：在 Minecraft 中学习技能，将成功的代码（如``制作剑''）存入技能库，后续任务可直接调用。
    \item \textbf{Generative Agents}：25 个模拟人类的 Agent 在小镇中生活，每个 Agent 维护一份事件日志（episodic memory）和从经验中提炼的语义记忆（semantic memory），这些记忆影响其后续行为。
    \item \textbf{Fine-tuning 作为长期记忆}：也可以直接 fine-tune 模型参数，将经验"写入"神经网络权重。
\end{itemize}

\subsection{KOALA：Agent 的统一抽象}

Shunyu Yao 提出了 KOALA 框架，用三个组件统一描述任意 Agent：

\begin{enumerate}
    \item \textbf{Memory}：信息存储的位置（上下文窗口 + 外部存储 + 模型参数）
    \item \textbf{Action Space}：Agent 可以执行的动作集合
    \item \textbf{Decision-making Procedure}：给定记忆和动作空间，如何选择下一步动作
\end{enumerate}

\subsection{本章小结}

长期记忆使 Agent 从``无状态''的一次性求解器进化为能够积累经验、持续改进的学习系统。记忆可以是文本日志、技能代码库，甚至模型参数本身。

